problem:  
Let $(M^4,g(t))$, $t\in[0,T)$, be a compact oriented 4‑manifold evolving by Ricci flow that develops a finite‑time singularity at $T<\infty$, and suppose $(M,g(t))$ is $\kappa$‑noncollapsed at all scales.  

Consider a sequence of blow‑ups around points $(x_i,t_i)\to(x_\ast,T)$ with scaling factors $\lambda_i = |{\rm Rm}(x_i,t_i)|$, producing pointed smooth limits  
\[
(N^4,g_\infty(t)),\quad t\in(-\infty,0],\qquad g_\infty(t)=\lim_{i\to\infty}\lambda_i\, g\bigl(t_i+\tfrac{t}{\lambda_i}\bigr).
\]
Assume each $(N,g_\infty(t))$ is nonflat, noncompact, and ancient, with $\mathrm{Rm}$ bounded.  

Define the **scale‑invariant curvature energy**
\[
\mathcal{E}(t)=\int_N |{\rm Rm}_{g_\infty(t)}|^2\, d\mu_{g_\infty(t)},
\]
which is invariant under rescaling in dimension four.  

**Problem (Finite‑energy rigidity conjecture):**  
Assume $\mathcal{E}(t)<\infty$ for some (and hence all) $t\le 0$.  
Must $(N,g_\infty(t))$ necessarily be a gradient shrinking soliton, up to scaling and time‑translation?  
Equivalently, can there exist a non‑solitonic ancient limit of 4‑dimensional Ricci flow whose total $L^2$ curvature energy is finite?  

Why is it a "good" problem:  
* **Natural invariant and analytic depth:** The total $L^2$ curvature is the *only* scale‑invariant quadratic curvature functional in dimension four, appearing in conformal geometry and the Chern–Gauss–Bonnet identity.  Asking whether finite $L^2$ curvature forces soliton structure is therefore a precise and analytically natural rigidity question at the critical dimension.  
* **Conceptual sharpness:** The problem gives a clean conjectural dichotomy—finite curvature energy ⇒ soliton—capturing the boundary between analytic rigidity and potential new exotic ancient solutions.  
* **Connection with known theory:** It builds directly on Perelman’s and subsequent rigidity results for κ‑solutions and shrinking solitons, but ventures beyond known Type‑I or monotonicity arguments, suggesting that *energy finiteness alone* might dictate soliton behavior.  
* **Bridge across disciplines:** It unites geometric analysis (Ricci flow blow‑up models) with four‑dimensional topology (through the $L^2$ curvature energy controlling Euler characteristic) and with nonlinear PDE rigidity methods (energy quantization and elliptic regularity).  
* **Simplicity and depth:** The question is stated in a single line yet is profoundly deep: “Is every finite‑energy nonflat ancient limit in dimension four a soliton?”  A positive resolution would show that Ricci flow singularities in 4D cannot produce exotic finite‑energy bubbles beyond shrinkers, a potential milestone unifying analytic and topological curvature quantization phenomena.